\usepackage{needspace}% used to keep the 10-rectangle working together
% =====================================================================
% preamble_local.tex: extra packages, styles and macros for the
% Definite Integrals chapter ONLY. Loaded after preamble_build.tex.
% Plot styles and the mini 3D macros are in figstyles.tex (shared with the
% slides and worksheets).
% =====================================================================
\usepgfplotslibrary{fillbetween}
\usepackage{adjustbox} % mini 3D pictures (figstyles.tex)
\usetikzlibrary{arrows.meta, decorations.pathreplacing, calc}
% =====================================================================
% figstyles.tex: plot styles and figure macros shared by the notes, the
% slides (../slides) and the worksheets (../worksheet). Needs pgfplots,
% the fillbetween library and tikz libraries arrows.meta and calc.
% House rules built in:
%   * `axis on top' in every sketch: shading never covers an axis, and the
%     fill starts at the axis line (items 31, 32, 34);
%   * regions are filled from the curve itself (fill between / closed plot);
%   * 3D solids in the page-13 style (translucent blue, light end caps).
% =====================================================================
\pgfplotsset{
  sketch/.style={clip=false, axis line style=thick, trig format=rad,
    axis lines=middle, axis on top,
    every axis x label/.style={at={(current axis.right of origin)},anchor=west},
    every axis y label/.style={at={(current axis.above origin)},anchor=south},
    ticks=none, xlabel=$x$, ylabel=$y$, width=7.6cm, height=6cm,
    every node/.append style={font=\footnotesize}},
  curve/.style={blue, semithick, smooth, samples=150},
  line2/.style={red!75!black, semithick, samples=2},
  % 3D solids of revolution (angles in DEGREES: no `trig format=rad' here)
  solid3d/.style={hide axis, trig format=deg, width=8.4cm, height=6.4cm, clip=false,
    z buffer=sort, colormap={bl}{rgb255=(150,185,235) rgb255=(40,90,180)},
    every node/.append style={font=\footnotesize}},
  solidsurf/.style={surf, shader=interp, opacity=0.75, point meta=y},
  red3d/.style={surf, shader=flat, draw=none, point meta=1,
    colormap={r}{rgb255=(215,50,50) rgb255=(215,50,50)}},
  orange3d/.style={surf, shader=interp, point meta=z,
    colormap={o}{rgb255=(250,210,160) rgb255=(220,120,30)}},
  cap3d/.style={surf, shader=flat, draw=none, point meta=1,
    colormap={c}{rgb255=(185,205,238) rgb255=(185,205,238)}},
  mini3d/.style={solid3d, width=2.6cm, height=2.4cm},
}
\tikzset{
  region/.style={fill=blue!22, draw=none},
  regionB/.style={fill=orange!35, draw=none},
  strip/.style={fill=red!55, draw=red!70!black, thin},
  asym/.style={black, dashed, thin},
  pt/.style={circle, fill=black, inner sep=1.1pt},
  dimarrow/.style={<->, >=Stealth, thin},
  ax3/.style={-{Stealth[length=2mm]}, thick},
}
\definecolor{hl}{RGB}{30,110,200}

% ---------------------------------------------------------------------
% Mini 3D solids for SOLUTIONS (items 54, 62, 66, 74): solid only, no
% washer, every picture 1.5 cm tall; a "big minus small" pair is 4.4 cm wide.
%  \solX[x unit cm]{R(x)}{a}{b}{R(a)}{R(b)}{x-axis end}{y-axis top}{curve colour}
%     solid swept out by y = R(x), a <= x <= b, about the x-axis
%  \solY[y unit cm]{R(y)}{c}{d}{R(c)}{R(d)}{x-axis end}{y-axis top}{curve colour}
%     solid swept out by x = R(y), c <= y <= d, about the y-axis
%     (R is written in terms of the letter x, which here stands for y)
% ---------------------------------------------------------------------
\newcommand{\solX}[9][0.5]{%
\begin{tikzpicture}[baseline=(current bounding box.center)]
\begin{axis}[mini3d, x={(#1cm,-0.06cm)}, y={(0.12cm,0.09cm)}, z={(0cm,0.4cm)}]
  \addplot3[cap3d, domain=0:1, y domain=0:360, samples=2, samples y=40] ({#3}, {x*(#5)*cos(y)}, {x*(#5)*sin(y)});
  \addplot3[solidsurf, opacity=0.5, point meta=z, domain=#3:#4, y domain=0:360, samples=24, samples y=40] ({x}, {(#2)*cos(y)}, {(#2)*sin(y)});
  \addplot3[cap3d, domain=0:1, y domain=0:360, samples=2, samples y=40] ({#4}, {x*(#6)*cos(y)}, {x*(#6)*sin(y)});
  \addplot3[#9, thick, domain=#3:#4, samples=40, samples y=1] ({x}, {0}, {#2});
  \draw[ax3] (axis cs:{min(0,#3)-0.3},0,0) -- (axis cs:#7,0,0);
  \draw[ax3] (axis cs:0,0,{-0.15*(#8)}) -- (axis cs:0,0,#8);
\end{axis}
\end{tikzpicture}}
\newcommand{\solY}[9][0.6]{%
\begin{tikzpicture}[baseline=(current bounding box.center)]
\begin{axis}[mini3d, x={(0.3cm,-0.03cm)}, y={(0.13cm,0.11cm)}, z={(0cm,#1cm)}]
  \addplot3[cap3d, domain=0:1, y domain=0:360, samples=2, samples y=40] ({x*(#5)*cos(y)}, {x*(#5)*sin(y)}, {#3});
  \addplot3[solidsurf, opacity=0.5, point meta=z, domain=#3:#4, y domain=0:360, samples=24, samples y=40] ({(#2)*cos(y)}, {(#2)*sin(y)}, {x});
  \addplot3[cap3d, domain=0:1, y domain=0:360, samples=2, samples y=40] ({x*(#6)*cos(y)}, {x*(#6)*sin(y)}, {#4});
  \addplot3[#9, thick, domain=#3:#4, samples=40, samples y=1] ({#2}, {0}, {x});
  \draw[ax3] (axis cs:0,0,{min(0,#3)-0.15*(#8)}) -- (axis cs:0,0,#8);
  \draw[ax3] (axis cs:-0.3,0,0) -- (axis cs:#7,0,0);
\end{axis}
\end{tikzpicture}}
\newcommand{\minisolo}[1]{\adjustbox{totalheight=1.5cm, max width=4.6cm}{#1}}
\newcommand{\minipair}[2]{\adjustbox{totalheight=1.5cm, max width=4.6cm}{#1\;\raisebox{0.15ex}{\Large$-$}\;#2}}
\newcommand{\minilabel}[1]{\raisebox{0.65cm}{\small #1}\;}


% ---------------------------------------------------------------------
% [limitsum] REVIEW-COPY SWITCH. DI_notes_REVIEW.tex defines \reviewcopy
% as 1 before loading this file. In the teacher and student copies the
% macros below do nothing (no highlight, no annotation).
%   \chg{note}{inline text}     : yellow highlight + hover note (one line only)
%   \begin{chgblock}{note} ... \end{chgblock}
%                               : light-yellow block + hover note covering it
%                                 + a sticky-note icon in the left margin
% The hover notes are PDF annotations written with xdvipdfmx \special's
% (no extra package needed).
% ---------------------------------------------------------------------
\providecommand\reviewcopy{0}
\ifnum\reviewcopy=1
  \tcbuselibrary{skins,breakable}
  \definecolor{chgyellow}{RGB}{255,246,170}
  \newcommand\chgann[3]{\special{pdf:ann width #1 height #2 depth 0pt
    << /Type/Annot /Subtype/Square /Contents (#3) /T (Review: limit-of-sum revision)
       /BS << /W 0 >> /F 4 >>}}
  \newcommand\chgicon[1]{\special{pdf:ann width 11pt height 11pt depth 0pt
    << /Type/Annot /Subtype/Text /Name/Comment /Contents (#1)
       /T (Review: limit-of-sum revision) /C [1 0.8 0] /F 4 >>}}
  \newbox\chgbox
  \newcommand\chg[2]{\leavevmode\begingroup\setbox\chgbox\hbox{#2}%
    \rlap{\color{chgyellow}\rule[\dimexpr-\dp\chgbox-1pt\relax]{\wd\chgbox}{\dimexpr\ht\chgbox+\dp\chgbox+2pt\relax}}%
    \chgann{\the\wd\chgbox}{\the\dimexpr\ht\chgbox+1pt\relax}{#1}\box\chgbox\endgroup}
  \newenvironment{chgblock}[1]{\par\begin{tcolorbox}[enhanced, breakable, colback=chgyellow!55,
      colframe=chgyellow!70!orange, boxrule=0.4pt, arc=0pt, boxsep=0pt,
      left=3pt, right=3pt, top=2pt, bottom=2pt, enlarge left by=-3pt, enlarge right by=-3pt,
      before skip=0pt, after skip=0pt, parbox=false, before upper={\parskip=0.9ex plus 0.3ex minus 0.2ex},
      overlay={\path let \p1=(frame.south west), \p2=(frame.north east) in
          node[anchor=south west, inner sep=0, outer sep=0] at (\p1)
          {\chgann{\the\dimexpr\x2-\x1\relax}{\the\dimexpr\y2-\y1\relax}{#1}};
        \node[anchor=north east, inner sep=0] at ([xshift=-30pt]frame.north west)
          {\chgicon{#1}\rule{0pt}{11pt}\hspace{11pt}};}]}%
    {\end{tcolorbox}\par}
\else
  \newcommand\chg[2]{#2}
  \newenvironment{chgblock}[1]{}{}
\fi
% \chgremoved{note}{marker text}: marks the spot where something was REMOVED
% (review copy only: small highlighted marker + hover note; prints nothing otherwise)
\ifnum\reviewcopy=1
  \newcommand\chgremoved[2]{\par\noindent\chg{#1}{\footnotesize\itshape [#2]}\par}
\else
  \newcommand\chgremoved[2]{}
\fi

% Small fractions are allowed only in exponents and in the limits of an
% integral (which are superscripts/subscripts); everywhere else use \dfrac.
\newcommand{\dx}{\,dx}
\newcommand{\dy}{\,dy}
\newcommand{\evalat}[3]{\left[\,#1\,\right]_{#2}^{#3}}

% [item 78] right-aligned A-level style marks at the end of a question part
% [limitsum] the mark allocation is now bold, e.g. [3]; in the review copy it
% is also highlighted with a hover note.
\newcommand{\mk}[1]{{\unskip\nobreak\hfil\penalty50\hskip1.5em\hbox{}\nobreak\hfill\chg{Mark allocation now bold}{\textbf{[#1]}}%
  \parfillskip=0pt\finalhyphendemerits=0\par}}
% [item 11] visually emphasised step labels in worked solutions
\newcommand{\step}[1]{\textcolor{teal!55!black}{\textbf{Step #1}}\quad}
% (a), (b), (c) part lists inside questions and solutions
\newenvironment{parts}{\begin{enumerate}[(a)]\setlength{\itemsep}{0.4ex}}{\end{enumerate}}
% SEAB-style solution layout (approved 10.5.8 template): working on the
% left, region sketch on the right (mini 3D pictures go under the sketch)
\newcommand{\solpair}[3][0.56]{\noindent\begin{minipage}[t]{#1\textwidth}\vspace{0pt}#2\end{minipage}\hfill
  \begin{minipage}[t]{\dimexpr0.98\textwidth-#1\textwidth\relax}\vspace{0pt}\centering #3\end{minipage}\par}
% minus-diagram style (drafts/minus_style_template.tex, items 29/31/62)
\pgfplotsset{minus/.style={axis lines=middle, width=5.0cm, height=4.2cm, ticks=none,
  xlabel={$x$}, ylabel={$y$}, every axis x label/.style={at={(ticklabel* cs:1)},anchor=west},
  every axis y label/.style={at={(ticklabel* cs:1)},anchor=south}, clip=false, axis on top,
  axis line style={gray!70}, every axis label/.append style={gray},
  every node/.append style={text=gray!80!black, font=\footnotesize}}}
\tikzset{hlreg/.style={fill=hl!35, draw=hl, thick}}

% ---------------------------------------------------------------------
% [limitsum] STUDENT COPY: blank working space sized to the teacher's
% solution. The hidden solution is typeset into a box that is never
% printed, only measured; the student copy then shows a light empty
% "Working" frame of about that height (at least \workspace, at most
% 0.8\textheight). The teacher copy is unchanged.
% ---------------------------------------------------------------------
\newsavebox\solmeasurebox
\newlength\solworkht
\newtcolorbox{workframe}[1]{enhanced, breakable=false, colback=white, colframe=black!22,
  boxrule=0.4pt, arc=1mm, left=2mm, right=2mm, top=1mm, bottom=1mm, height=#1,
  before skip=1mm, after skip=3mm}
\ifx\conform\ShowSolutionsKey\else
  \renewcommand\solution[1]{\par
    \begingroup
      \setbox\solmeasurebox\vbox{\hsize\linewidth\linewidth\hsize
        \noindent\textbf{Solution.}\par\vspace{0.5ex}\noindent #1\par}%
      \global\solworkht\dimexpr\ht\solmeasurebox+\dp\solmeasurebox+4mm\relax
    \endgroup
    \ifdim\solworkht<\workspace \global\solworkht\workspace\fi
    \ifdim\solworkht>0.8\textheight \global\solworkht0.8\textheight\fi
    \begin{workframe}{\solworkht}{\small\color{black!45}\textbf{Working}}\end{workframe}\par}
\fi
